\chapter{CƠ SỞ LÝ THUYẾT}
\label{chap:2-background}

Chương này trình bày nền tảng lý thuyết của kỹ thuật xử lý micro-batch, là kỹ thuật xử lý dữ liệu mà đề tài áp dụng cho lĩnh vực nông nghiệp. Mục~\ref{sec:2-challenges} nêu các thách thức của xử lý luồng thuần túy, là động lực dẫn tới micro-batch. Mục~\ref{sec:2-definition} định nghĩa micro-batch. Mục~\ref{sec:2-mechanism} giải thích cơ chế hoạt động, từ kiến trúc, vòng đời một lô đến cách xử lý lỗi và dữ liệu đến muộn. Mục~\ref{sec:2-tech} giới thiệu bộ công nghệ Apache Kafka và Spark Structured Streaming. Mục~\ref{sec:2-compare} so sánh micro-batch với xử lý theo lô và xử lý luồng, cùng các đánh đổi. Mục~\ref{sec:2-dbt-quality} trình bày dbt và kỹ thuật bảo đảm chất lượng dữ liệu. Mục~\ref{sec:2-lakehouse} trình bày lakehouse gồm Apache Iceberg, HDFS, Hive Metastore và Trino, là nơi lưu lịch sử và phục vụ phân tích.

\section{Thách thức của xử lý luồng thuần túy}
\label{sec:2-challenges}

Trong xử lý luồng thuần túy (pure stream processing), mỗi sự kiện được xử lý ngay khi đến. Cách làm này cho độ trễ thấp, nhưng khi áp dụng cho dữ liệu cảm biến nông nghiệp thì nảy sinh ba khó khăn chính.

\subsection{Khôi phục trạng thái đang xử lý}
Engine luồng giữ trạng thái của từng khóa (bộ đếm, cửa sổ, giá trị gần nhất) trong bộ nhớ của các nút xử lý. Để khôi phục sau sự cố, engine phải chụp trạng thái này định kỳ, và khi khởi động lại phải nạp lại toàn bộ trạng thái rồi phát lại các sự kiện đến sau lần chụp cuối. Trạng thái càng lớn thì việc chụp và nạp lại càng tốn thời gian. Ví dụ, khi một cổng LoRaWAN gửi dồn dữ liệu của nhiều giờ mất kết nối, trạng thái của mọi thiết bị liên quan phải được giữ và chụp lại trong lúc engine xử lý đợt dữ liệu đó.

\subsection{Ghi trùng khi ghi dở dang}
Khi xử lý từng sự kiện, không có ranh giới tự nhiên cho biết phần nào đã được ghi ra đích. Nếu tác vụ sập sau khi ghi một sự kiện nhưng trước khi lưu tiến độ, sự kiện đó sẽ được xử lý và ghi lại sau khi khởi động lại. Muốn tránh, đích phải tham gia giao thức cam kết hai pha cùng checkpoint của engine, hoặc mọi lệnh ghi phải idempotent. Ví dụ, nếu nhật ký lượng nước tưới được cập nhật bằng phép cộng dồn, một sự kiện được ghi lại sẽ bị cộng hai lần.

\subsection{Chi phí của độ trễ dưới một giây}
Để đạt độ trễ dưới một giây, engine phải chạy liên tục, giữ trạng thái trong bộ nhớ và xử lý dữ liệu đến sai thứ tự bằng watermark cho từng sự kiện. Chi phí này không đáng với những tác vụ không cần độ trễ đó. Ví dụ, nhiệt độ đất biến thiên chậm, nên chu kỳ cập nhật 30 giây đã đáp ứng nhu cầu canh tác.

\section{Khái niệm xử lý micro-batch}
\label{sec:2-definition}

Micro-batch là mô hình xử lý lai: chia một luồng dữ liệu liên tục, không giới hạn thành chuỗi các khối dữ liệu nhỏ, hữu hạn (micro-batch), mỗi khối được thực thi sau một khoảng thời gian ngắn. Kleppmann mô tả ý tưởng này là ``chia luồng thành các khối nhỏ và xử lý mỗi khối như một tiến trình xử lý theo lô thu nhỏ'' \cite{kleppmann2017ddia}.

Nhờ ranh giới rõ ràng của từng khối, mỗi lô có thể được ghi nhận, thử lại và cam kết (commit) như một đơn vị độc lập. Mô hình kết hợp độ trễ thấp của xử lý luồng với thông lượng cao và khả năng chịu lỗi của xử lý theo lô \cite{conduktor_microbatch}.

\begin{figure}[H]
\centering
\hinhtrong{0.85\linewidth}{4.5cm}{Luồng sự kiện liên tục được chia thành các micro-batch theo chu kỳ trigger}
\caption{Rời rạc hóa luồng dữ liệu thành các micro-batch}
\label{fig:2-microbatch-concept}
\end{figure}

\section{Cơ chế hoạt động}
\label{sec:2-mechanism}

\subsection{Kiến trúc pipeline}
\label{sec:2-pipeline}

Một pipeline micro-batch điển hình cho dữ liệu cảm biến gồm các thành phần nối tiếp trong Hình~\ref{fig:2-pipeline}. Dữ liệu từ cảm biến qua cổng kết nối (gateway) vào nhật ký sự kiện Kafka. Kafka tiếp tục nhận sự kiện trong khi Spark đang tính toán, nhờ đó khâu thu thập và khâu xử lý tách rời nhau. Kết quả được ghi vào cơ sở dữ liệu phục vụ bảng điều khiển. Tiến độ xử lý và trạng thái được lưu bền vững trong checkpoint.

\begin{figure}[H]
\centering
\begin{tikzpicture}[
  node distance=0.55cm,
  box/.style={rectangle, rounded corners, draw, minimum width=2.0cm, minimum height=1.0cm, align=center, font=\small, fill=gray!10},
  arr/.style={-{Stealth}, thick}]
\node[box] (s) {Cảm biến};
\node[box, right=of s] (g) {Gateway};
\node[box, right=of g, fill=blue!10] (k) {Kafka\\nhật ký sự kiện};
\node[box, right=of k, fill=orange!15] (p) {Spark\\tính toán};
\node[box, right=of p] (d) {CSDL \&\\dashboard};
\node[box, below=0.9cm of p, fill=green!10] (c) {Checkpoint\\offset + trạng thái};
\draw[arr] (s)--(g); \draw[arr] (g)--(k); \draw[arr] (k)--(p); \draw[arr] (p)--(d);
\draw[arr, dashed] (p)--(c);
\end{tikzpicture}
\caption{Kiến trúc pipeline micro-batch từ cảm biến đến bảng điều khiển}
\label{fig:2-pipeline}
\end{figure}

\subsection{Vòng đời xử lý của một micro-batch}
\label{sec:2-lifecycle}

Mỗi micro-batch trải qua sáu bước:
\begin{enumerate}
  \item \textbf{Trigger}: bắt đầu lượt chạy tiếp theo.
  \item \textbf{Lập kế hoạch và ghi offset}: chọn các bản ghi mới và ghi lại phạm vi offset sẽ đọc.
  \item \textbf{Đọc và biến đổi}: áp dụng truy vấn lên dữ liệu của lô.
  \item \textbf{Cập nhật trạng thái}: dành cho các truy vấn có trạng thái (cửa sổ, tổng hợp, khử trùng lặp).
  \item \textbf{Ghi kết quả}: gửi đầu ra tới nơi lưu (sink).
  \item \textbf{Cam kết lô}: ghi nhận lô đã hoàn tất.
\end{enumerate}

\subsection{Chu kỳ trigger và ví dụ dữ liệu nông nghiệp}
\label{sec:2-trigger}

Chu kỳ trigger quyết định độ chi tiết của lô. Giả sử một cảm biến độ ẩm đất gửi số đo lên Kafka, và trigger chạy mỗi 5 giây (Bảng~\ref{tab:2-example}). Mỗi lô tính một giá trị tổng hợp rồi lưu kết quả, hiển thị trên bảng điều khiển và đưa vào các quy tắc cảnh báo.

\begin{table}[H]
\centering
\caption{Ví dụ hai micro-batch 5 giây của một cảm biến độ ẩm đất}
\label{tab:2-example}
\small
\begin{tblr}{
  colspec = {X[1,l]X[1.3,l]X[2.2,l]X[1.5,l]}, rows = {abovesep=2pt, belowsep=2pt},
  hlines, vlines, row{1} = {font=\bfseries, bg=gray!20},
}
Lô & Trigger & Số đo đến Kafka (độ ẩm \%) & Trung bình \\
1 & 08:00:05 & 22 (08:00:01), 24 (08:00:03), 26 (08:00:04) & 24\% \\
2 & 08:00:10 & 28 (08:00:07), 30 (08:00:09) & 29\% \\
\end{tblr}
\end{table}

\subsection{Độ tin cậy: phục hồi sau lỗi và dữ liệu đến muộn}
\label{sec:2-reliability}

\paragraph{Phục hồi sau lỗi.}
Checkpoint bền vững lưu offset và trạng thái truy vấn. Khi khởi động lại, hệ thống phát lại và hoàn tất lô chưa được cam kết, nên không mất và không bỏ sót dữ liệu.

\paragraph{Dữ liệu đến muộn.}
Cần phân biệt thời gian sự kiện (event time) với thời gian đến (arrival time). Một số đo có thể được tạo lúc 08:00:03 nhưng đến hệ thống sau đó vì mất kết nối. Hệ thống giữ trạng thái cửa sổ để nhận các số đo muộn này. \emph{Watermark} đi theo tiến độ của thời gian sự kiện; nó không phải bộ đếm chờ theo đồng hồ thực. Khi watermark tiến lên, trạng thái của các cửa sổ đủ điều kiện được dọn, giữ bộ nhớ ở mức giới hạn.

\paragraph{Điều kiện exactly-once.}
Kết quả ``đúng một lần'' (exactly-once) cần đồng thời ba điều kiện: nguồn có thể phát lại (replayable), engine có checkpoint, và nơi ghi an toàn khi thử lại (idempotent hoặc transactional). Nếu chỉ dùng sink Kafka thông thường thì chỉ đạt at-least-once.

\section{Công nghệ hiện thực}
\label{sec:2-tech}

Bộ ba công nghệ được dùng để hiện thực micro-batch gồm Apache Kafka làm nguồn, Spark Structured Streaming làm engine xử lý và checkpoint bảo đảm độ tin cậy.

\subsection{Apache Kafka làm nguồn dữ liệu}
\label{sec:2-kafka}

Kafka là nhật ký sự kiện bền vững và có thể phát lại \cite{kafka_docs}. Dữ liệu được tổ chức theo \emph{topic}, mỗi topic chia thành các \emph{partition}; mỗi thông điệp trong partition có một \emph{offset} tăng dần. Thông điệp được giữ trong một thời gian lưu trữ (retention) kể cả sau khi đã đọc, nên Spark có thể đọc lại bất kỳ khoảng offset nào. Kafka cũng tách bên sinh dữ liệu khỏi bên xử lý. Một partition Kafka tương ứng một tác vụ (task) Spark.

Ở đầu mỗi lô, Spark hỏi Kafka offset mới nhất của từng partition, rồi xác định khoảng offset cần đọc (ví dụ partition 0: 1000--1500, partition 1: 980--1420). Khoảng offset này chính là micro-batch. Các partition được đọc song song, và offset được theo dõi bằng checkpoint của Spark chứ không dùng consumer group. Tùy chọn \code{maxOffsetsPerTrigger} giới hạn số bản ghi mỗi lô, bảo vệ tác vụ khỏi các đợt tăng đột biến.

\subsection{Spark Structured Streaming}
\label{sec:2-spark}

Spark Structured Streaming coi luồng là một \emph{bảng không giới hạn}: mỗi bản ghi mới là một dòng được nối vào bảng đầu vào. Ở mỗi trigger, Spark chạy truy vấn chỉ trên các dòng mới và cập nhật kết quả \cite{spark_ss}. Nhờ đó cùng một API DataFrame/SQL dùng được cho cả xử lý theo lô và xử lý luồng. Ba khái niệm cấu hình chính:

\begin{itemize}
  \item \textbf{Trigger} (khi nào chạy): \code{processingTime} chạy theo chu kỳ cố định, ví dụ mỗi 10 giây; mặc định là các lô nối tiếp nhau; \code{availableNow} xử lý hết dữ liệu hiện có rồi dừng.
  \item \textbf{Output mode} (ghi gì): \emph{append} chỉ ghi dòng mới, \emph{update} chỉ ghi dòng thay đổi, \emph{complete} ghi lại toàn bộ bảng kết quả.
  \item \textbf{Watermark} (xử lý dữ liệu muộn): ví dụ ``chờ tối đa 2 phút rồi chốt cửa sổ''; trạng thái cũ bị xóa nên bộ nhớ không tăng vô hạn.
\end{itemize}

\subsection{Checkpoint}
\label{sec:2-checkpoint}

Thư mục checkpoint lưu bốn thành phần:
\begin{itemize}
  \item \code{offsets/}: nhật ký ghi trước (write-ahead log). \emph{Trước} khi lô $N$ chạy, Spark ghi phạm vi offset sẽ đọc.
  \item \code{commits/}: \emph{sau} khi lô $N$ được ghi vào sink, Spark đánh dấu lô đã xong.
  \item \code{state/}: kho trạng thái cho cửa sổ, tổng hợp và khử trùng lặp.
  \item \code{metadata}: mã định danh duy nhất của truy vấn. Khi đổi lớn logic truy vấn, cần dùng vị trí checkpoint mới.
\end{itemize}

Quy trình phục hồi như sau: ghi offset vào \code{offsets/}, xử lý và ghi vào sink, rồi ghi \code{commits/}. Nếu hệ thống sập khi lô $N$ đã có trong \code{offsets/} nhưng chưa có trong \code{commits/}, lô $N$ được phát lại với \emph{cùng} khoảng offset. Để kết quả đúng một lần, sink phải idempotent hoặc transactional, ví dụ một bảng Iceberg hay Delta Lake ghi theo giao dịch, hoặc lệnh upsert theo khóa.

\subsection{Ví dụ mã nguồn}
\label{sec:2-code}

Đoạn mã dưới đây đọc luồng đơn hàng từ Kafka, tính doanh thu theo sản phẩm trong cửa sổ 1 phút và cập nhật kết quả mỗi 10 giây. Tác vụ của đề tài (Chương~\ref{chap:4-design}) dùng cùng nguồn Kafka, \code{maxOffsetsPerTrigger} và trigger \code{processingTime}, nhưng không tổng hợp theo cửa sổ: mỗi lô được xử lý trong \code{foreachBatch} và ghi bằng lệnh \code{MERGE}.

\begin{lstlisting}[language=Python, caption={Kafka $\rightarrow$ Spark $\rightarrow$ sink với trigger 10 giây}, label={lst:2-example}]
orders = (spark.readStream.format("kafka")
    .option("subscribe", "orders")
    .option("maxOffsetsPerTrigger", 10000)
    .load())

revenue = (orders
    .select(from_json("value", schema).alias("o")).select("o.*")
    .withWatermark("event_time", "2 minutes")
    .groupBy(window("event_time", "1 minute"), "product")
    .agg(sum("amount")))

(revenue.writeStream
    .outputMode("update")
    .trigger(processingTime="10 seconds")
    .option("checkpointLocation", "/tmp/chk/orders")
    .start())
\end{lstlisting}

\section{So sánh, trường hợp sử dụng và đánh đổi}
\label{sec:2-compare}

\subsection{So sánh ba mô hình xử lý}
\label{sec:2-three-models}

\begin{table}[H]
\centering
\caption{So sánh xử lý theo lô, micro-batch và xử lý luồng}
\label{tab:2-compare}
\small
\begin{tblr}{
  colspec = {X[1.2,l]X[2,l]X[2,l]X[2,l]}, rows = {abovesep=2pt, belowsep=2pt},
  hlines, vlines, row{1} = {font=\bfseries, bg=gray!20},
}
Tiêu chí & Xử lý theo lô & Micro-batch & Xử lý luồng \\
Đơn vị xử lý & Khối dữ liệu lớn & Lô nhỏ, thường xuyên & Từng sự kiện liên tục \\
Độ trễ & Phút đến giờ & Giây đến phút & Mili giây đến giây \\
Thông lượng & Cao, hiệu quả & Tốt, độ phức tạp vừa phải & Phụ thuộc thiết kế \\
Độ phức tạp & Đơn giản & Cân bằng độ trễ và hiệu quả & Cao hơn, chi phí vận hành lớn \\
Phù hợp với & ETL, báo cáo, phân tích lịch sử & Dashboard, ETL tăng dần, phân tích gần thời gian thực & Phát hiện gian lận, cảnh báo, quyết định thời gian thực \\
\end{tblr}
\end{table}

\subsection{Trường hợp sử dụng}
\label{sec:2-usecases}
Micro-batch thường dùng cho: phân tích gần thời gian thực, xử lý nhật ký, giám sát IoT, ETL tăng dần vào kho dữ liệu, và các pipeline học máy.

\subsection{Được và mất}
\label{sec:2-tradeoffs}

\begin{table}[H]
\centering
\caption{Lợi ích và chi phí của micro-batch}
\label{tab:2-tradeoffs}
\small
\begin{tblr}{
  colspec = {X[1,l]X[1,l]}, rows = {abovesep=2pt, belowsep=2pt},
  hlines, vlines, row{1} = {font=\bfseries, bg=gray!20},
}
Lợi ích & Chi phí \\
Dữ liệu mới trong vài phút với công cụ theo lô quen thuộc & Độ trễ không thể thấp hơn chu kỳ lô \\
Thử lại dễ dàng, cam kết idempotent và transactional & Cam kết thường xuyên tạo nhiều tệp nhỏ và snapshot, cần gộp tệp (compaction) và dọn snapshot \\
Rẻ và đơn giản hơn xử lý luồng thực sự & Chu kỳ càng ngắn thì chi phí lập lịch và siêu dữ liệu càng cao \\
\end{tblr}
\end{table}

\subsection{Liên hệ với đề tài}
\label{sec:2-link}
Dữ liệu nông nghiệp như độ ẩm đất hay nhiệt độ biến thiên chậm, kết nối LoRaWAN dễ gián đoạn, và yêu cầu nghiệp vụ thường chấp nhận độ trễ từ vài giây đến vài chục giây. Các đặc điểm này phù hợp với micro-batch: tránh chi phí của xử lý luồng thuần túy mà vẫn bảo đảm tính liên tục và đúng đắn của dữ liệu. Việc áp dụng cụ thể được trình bày ở Chương~\ref{chap:4-design}.

\section{dbt và kỹ thuật bảo đảm chất lượng dữ liệu}
\label{sec:2-dbt-quality}

Phần trên giải thích cách dữ liệu được đưa vào hệ thống theo từng micro-batch. Mục này trình bày nhóm kỹ thuật bảo đảm dữ liệu đó \emph{đáng tin cậy}: kiến trúc phân tầng medallion, công cụ biến đổi dbt, cách kiểm định, cách ly bản ghi lỗi và lưu vết nguồn gốc (lineage). Đây là cơ sở cho yêu cầu BR4 và cho phần đánh giá tính đúng đắn ở Chương~\ref{chap:5-evaluation}.

\subsection{Kiến trúc medallion}
\label{sec:2-medallion}

Kiến trúc medallion chia kho dữ liệu thành ba tầng, mỗi tầng có một cam kết chất lượng rõ ràng \cite{databricks_medallion}:

\begin{itemize}
  \item \textbf{Bronze}: dữ liệu thô gần với nguồn, chỉ bổ sung siêu dữ liệu (thời điểm nhận, nguồn, cờ chất lượng). Tầng này giữ toàn bộ lịch sử để có thể xử lý lại.
  \item \textbf{Silver}: dữ liệu đã làm sạch, khử trùng lặp và chuẩn hóa kiểu. Mỗi sự kiện chỉ xuất hiện một lần.
  \item \textbf{Gold}: dữ liệu tổng hợp theo nhu cầu nghiệp vụ (theo ngày, theo thửa ruộng), sẵn sàng cho bảng điều khiển và phân tích.
\end{itemize}

Ưu điểm chính là \emph{tách bạch trách nhiệm}: lỗi ở tầng sau luôn có thể sửa bằng cách tính lại từ tầng trước, vì tầng bronze không bao giờ bị ghi đè bằng dữ liệu đã biến đổi. Trong đề tài, bronze được ghi bởi tác vụ Spark Structured Streaming, còn silver và gold được dựng bằng dbt.

\subsection{dbt: biến đổi dữ liệu dưới dạng mã}
\label{sec:2-dbt}

dbt (data build tool) cho phép viết mỗi bước biến đổi thành một \emph{model}, tức một câu lệnh \code{SELECT} lưu trong tệp \code{.sql} \cite{dbt_docs}. dbt tự sinh câu lệnh tạo bảng, tự xác định thứ tự chạy nhờ hàm \code{ref()} và \code{source()}, và chạy câu lệnh trên một engine có sẵn. Đề tài dùng bộ chuyển đổi dbt-spark kết nối tới Spark Thrift Server, nên Spark vẫn là engine thực thi và các bảng đích là bảng Iceberg.

Đoạn mã~\ref{lst:2-silver} là model silver của đề tài. Hàm cửa sổ \code{row\_number()} giữ lại đúng một bản ghi cho mỗi khóa \code{(device\_id, metric, measured\_at)}, ưu tiên bản ghi được nhận sau cùng.

\begin{lstlisting}[language=SQL, caption={Model silver khử trùng lặp (\texttt{dbt/models/silver/sensor\_readings.sql}, rút gọn)}, label={lst:2-silver}]
{{ config(materialized='table') }}
with ranked as (
    select *,
        row_number() over (
            partition by device_id, metric, measured_at
            order by ingested_at desc, source_offset desc
        ) as row_num
    from {{ source('bronze', 'sensor_readings') }}
    where measured_at is not null
      and device_id is not null and metric is not null
)
select event_id, device_id, field_id, farm_id, crop_code, metric,
       value, unit, measured_at, ingested_at, is_late, quality_flag,
       source_topic, source_partition, source_offset, batch_id
from ranked
where row_num = 1
\end{lstlisting}

\subsection{Kiểm thử dữ liệu trong dbt}
\label{sec:2-dbt-tests}

Mỗi model khai báo các \emph{test} trong tệp \code{schema.yml}. Một test là một câu truy vấn trả về các dòng vi phạm; test đạt khi kết quả rỗng. dbt cung cấp sẵn bốn test chung: \code{not\_null}, \code{unique}, \code{accepted\_values} và \code{relationships}. Đoạn mã~\ref{lst:2-silver-tests} là phần khai báo test của model silver.

\begin{lstlisting}[caption={Test của model silver (\texttt{dbt/models/silver/schema.yml}, rút gọn)}, label={lst:2-silver-tests}]
models:
  - name: sensor_readings
    tests:
      - unique:
          column_name: "concat(device_id, '|', metric, '|', cast(measured_at as string))"
    columns:
      - name: device_id
        tests: [not_null]
      - name: metric
        tests:
          - not_null
          - accepted_values:
              values: [soil_moisture, soil_temperature, air_temperature,
                       air_humidity, relative_humidity, rainfall]
\end{lstlisting}

Lệnh \code{dbt build} chạy model và test theo đúng thứ tự phụ thuộc. Khi một test của model silver thất bại, các model gold phía sau bị bỏ qua (\emph{skip}) thay vì được dựng từ dữ liệu sai. Đây là cơ chế đề tài dùng để đáp ứng FR4.7: một test thất bại sẽ chặn lần làm mới tầng gold tương ứng.

\subsection{Kiểm định, cách ly và gắn cờ}
\label{sec:2-validation}

Không phải lỗi nào cũng nên xử lý giống nhau. Đề tài phân biệt hai mức:

\begin{itemize}
  \item \textbf{Lỗi cấu trúc} (payload không phải JSON, thiếu trường bắt buộc, thời gian không đọc được, giá trị không phải số, thiết bị không có trong danh mục): bản ghi không thể gán cho thửa ruộng nào, nên được \emph{cách ly} (quarantine) sang một bảng riêng kèm lý do và payload gốc. Tác vụ không dừng lại vì một bản ghi hỏng.
  \item \textbf{Lỗi nội dung} (giá trị rỗng, ngoài khoảng vật lý, bước nhảy bất thường, bản ghi trùng): bản ghi vẫn được lưu nhưng mang \emph{cờ chất lượng} (\code{quality\_flag}), để bảng điều khiển có thể loại bỏ hoặc đánh dấu.
\end{itemize}

Cách làm này giữ được toàn bộ bằng chứng: không có bản ghi nào bị xóa âm thầm, và số lượng của từng loại lỗi có thể đếm lại. Đoạn mã~\ref{lst:2-quarantine-ddl} là cấu trúc bảng cách ly.

\begin{lstlisting}[language=SQL, caption={Bảng cách ly trên Iceberg (\texttt{src/streaming/ddl.py}, rút gọn)}, label={lst:2-quarantine-ddl}]
CREATE TABLE IF NOT EXISTS lakehouse.ops.quarantine (
    reason string NOT NULL,      -- invalid_json | missing_field | ...
    raw_payload string,          -- original payload, unchanged
    device_id string, metric string, reading_ts string, value_raw string,
    source_topic string NOT NULL, source_partition int NOT NULL,
    source_offset bigint NOT NULL,
    ingested_at timestamp, detected_at timestamp, batch_id bigint
) USING iceberg PARTITIONED BY (days(detected_at))
\end{lstlisting}

\subsection{Lưu vết nguồn gốc (lineage)}
\label{sec:2-lineage}

Lineage là khả năng trả lời câu hỏi ``con số này đến từ đâu''. Ở mức bản ghi, đề tài gắn vào mọi dòng các cột sau: chủ đề, phân vùng và offset Kafka (\code{source\_topic}, \code{source\_partition}, \code{source\_offset}), thời điểm Kafka nhận sự kiện (\code{ingested\_at}), mã micro-batch đã ghi dòng đó (\code{batch\_id}) và cờ chất lượng. Vì offset Kafka là duy nhất trong một phân vùng, bộ ba (topic, partition, offset) định danh chính xác một sự kiện gốc. Kết hợp với khả năng phát lại của Kafka (Mục~\ref{sec:2-kafka}), mọi giá trị trên bảng điều khiển đều truy ngược được tới sự kiện và lô đã tạo ra nó.

\subsection{Ghi idempotent và tính đúng đắn khi phát lại}
\label{sec:2-idempotent}

Mục~\ref{sec:2-checkpoint} cho thấy một lô có thể được chạy lại sau sự cố. Để kết quả không bị nhân đôi, mọi lệnh ghi phải \emph{idempotent}: chạy một lần hay nhiều lần cho cùng kết quả. Đề tài dùng lệnh \code{MERGE INTO} của Iceberg với khóa tự nhiên của từng bảng, như trong Đoạn mã~\ref{lst:2-merge}.

\begin{lstlisting}[language=SQL, caption={Ghi idempotent bằng MERGE (\texttt{src/streaming/ddl.py})}, label={lst:2-merge}]
MERGE INTO lakehouse.bronze.sensor_readings AS t
USING src_bronze AS s
ON t.device_id = s.device_id AND t.metric = s.metric
   AND t.measured_at = s.measured_at
WHEN MATCHED THEN UPDATE SET *
WHEN NOT MATCHED THEN INSERT *
\end{lstlisting}

Khóa của bảng cách ly là (topic, partition, offset); khóa của bảng đếm chất lượng là (\code{batch\_id}, \code{device\_id}, \code{metric}). Cột \code{detected\_at} trong bảng cách ly lấy bằng \code{ingested\_at} chứ không lấy giờ hệ thống, để lần phát lại sinh ra dòng giống hệt lần đầu.

% Mục 2.7 (người phụ trách: Khang).
\section{Lakehouse: Apache Iceberg, HDFS, Hive Metastore và Trino}
\label{sec:2-lakehouse}

Các mục trên giải thích dữ liệu đi vào hệ thống và được làm sạch ra sao. Mục này trình bày nơi dữ liệu được \emph{lưu lâu dài} và cách nó được \emph{truy vấn} để phục vụ phân tích lịch sử, tức là cơ sở của BR5 và BR6. Bốn thành phần phối hợp với nhau: HDFS giữ tệp, Apache Iceberg tổ chức các tệp đó thành bảng, Hive Metastore ghi nhớ bảng nào nằm ở đâu, và Trino là engine truy vấn trên các bảng đó.

\subsection{Từ data lake đến lakehouse}
\label{sec:2-lake-to-lakehouse}

Một \emph{data lake} lưu dữ liệu dưới dạng tệp (thường là Parquet) trên hệ thống tệp phân tán. Cách này rẻ và co giãn tốt, nhưng thiếu những gì cơ sở dữ liệu quan hệ có sẵn: giao dịch, ghi đè an toàn, sửa schema và xem lại quá khứ. Với đường xử lý micro-batch, thiếu sót này có hậu quả cụ thể. Mỗi 30 giây một lô mới ghi thêm tệp; nếu người đọc mở bảng đúng lúc lô đang ghi dở, họ có thể thấy dữ liệu không nhất quán, và việc ghi lại một lô sau sự cố có thể sinh bản ghi trùng.

\emph{Lakehouse} giải quyết điều đó bằng cách đặt một \emph{định dạng bảng} (table format) lên trên tệp Parquet. Định dạng bảng là một lớp siêu dữ liệu cho biết \emph{tập tệp nào hiện là bảng}. Đề tài chọn Apache Iceberg \cite{iceberg_spec}; các lựa chọn tương tự là Delta Lake và Apache Hudi.

\subsection{Apache Iceberg}
\label{sec:2-iceberg}

\paragraph{Cấu trúc siêu dữ liệu.} Iceberg không liệt kê thư mục để biết bảng có tệp nào. Thay vào đó, một bảng là một cây siêu dữ liệu: tệp \emph{metadata} chứa schema, cách phân vùng và danh sách các \emph{snapshot}; mỗi snapshot trỏ tới một \emph{manifest list}; manifest list trỏ tới các \emph{manifest}; mỗi manifest liệt kê các tệp dữ liệu cùng thống kê cột (giá trị nhỏ nhất, lớn nhất, số dòng). Mỗi lần ghi (commit) tạo một snapshot mới và hoán đổi con trỏ tới tệp metadata theo kiểu nguyên tử.

\begin{figure}[H]
\centering
\begin{tikzpicture}[
  node distance=0.5cm and 0.55cm,
  box/.style={rectangle, rounded corners, draw, minimum width=2.3cm, minimum height=0.9cm, align=center, font=\footnotesize, fill=gray!10},
  arr/.style={-{Stealth}, thick}]
\node[box, fill=blue!10] (cat) {Hive Metastore\\con trỏ tới metadata};
\node[box, right=of cat] (meta) {Tệp metadata\\schema, phân vùng,\\danh sách snapshot};
\node[box, right=of meta] (ml) {Manifest list\\của snapshot};
\node[box, below=of ml] (mf) {Manifest\\tệp dữ liệu + thống kê};
\node[box, fill=green!10, left=of mf] (df) {Tệp Parquet\\trên HDFS};
\draw[arr] (cat)--(meta); \draw[arr] (meta)--(ml); \draw[arr] (ml)--(mf); \draw[arr] (mf)--(df);
\end{tikzpicture}
\caption{Cây siêu dữ liệu của một bảng Iceberg}
\label{fig:2-iceberg-tree}
\end{figure}

Cấu trúc này mang lại các khả năng mà thư mục Parquet thuần không có, và mỗi khả năng gắn với một yêu cầu của đề tài:

\begin{itemize}
  \item \textbf{Ghi nguyên tử và cô lập snapshot.} Người đọc luôn thấy một snapshot hoàn chỉnh. Lô streaming đang ghi dở không bao giờ lộ ra ngoài. Điều này cần cho BR3 và BR4.
  \item \textbf{Phân vùng ẩn (hidden partitioning).} Bảng được phân vùng theo một \emph{phép biến đổi} của cột, ví dụ \code{days(measured\_at)}. Người viết truy vấn chỉ lọc theo \code{measured\_at}, Iceberg tự bỏ qua các phân vùng không liên quan. Đây là cơ sở của FR6.2.
  \item \textbf{Tiến hóa schema.} Thêm, đổi tên hoặc bỏ cột chỉ sửa siêu dữ liệu, không viết lại dữ liệu cũ. Cột được định danh bằng số chứ không bằng tên.
  \item \textbf{Du hành thời gian (time travel).} Vì các snapshot cũ còn nguyên, có thể truy vấn bảng ở một thời điểm hoặc phiên bản trong quá khứ. Đây là cơ sở của việc tái tạo một báo cáo cũ (BR6).
  \item \textbf{Cập nhật theo dòng.} Phiên bản bảng 2 hỗ trợ xóa và cập nhật theo dòng, nên lệnh \code{MERGE INTO} dùng được cho ghi idempotent (Mục~\ref{sec:2-idempotent}).
\end{itemize}

\paragraph{Vấn đề tệp nhỏ và compaction.} Mỗi micro-batch là một commit, và một commit thường sinh ra ít nhất một tệp mới trong mỗi phân vùng nó chạm tới. Với trigger 30 giây, một ngày có khoảng 2\,880 commit; dữ liệu của cả ngày bị chia thành hàng nghìn tệp rất nhỏ. Truy vấn phải mở và đọc siêu dữ liệu của từng tệp nên chậm đi. Iceberg cung cấp thủ tục \code{rewrite\_data\_files} để gộp các tệp nhỏ thành tệp lớn (chiến lược \emph{binpack}) \cite{iceberg_procedures}. Việc gộp tạo một snapshot mới nên không làm gián đoạn người đọc đang truy vấn snapshot cũ. Hiệu quả của bước này được đo ở Mục~\ref{sec:5-bench-compaction}.

\subsection{HDFS làm kho lưu trữ}
\label{sec:2-hdfs}

HDFS (Hadoop Distributed File System) chia tệp thành các khối lớn (mặc định 128 MB), sao chép mỗi khối trên nhiều DataNode và để NameNode giữ danh mục tệp \cite{hdfs_design}. HDFS được tối ưu cho tệp lớn được ghi một lần và đọc tuần tự, nên \emph{không hợp} với hàng nghìn tệp nhỏ: mỗi tệp, dù nhỏ đến đâu, chiếm một mục trong bộ nhớ của NameNode. Đây là một lý do nữa để compaction là bắt buộc chứ không chỉ là tối ưu. Trong đề tài, một NameNode và một DataNode chạy trong Docker; hệ số sao chép là 1 vì chỉ có một nút.

\subsection{Hive Metastore làm catalog}
\label{sec:2-hms}

Iceberg cần một \emph{catalog} để trả lời câu hỏi ``bảng \code{gold.field\_daily} hiện có tệp metadata nào''. Hive Metastore (HMS) là một dịch vụ lưu danh mục bảng và con trỏ tới metadata hiện tại \cite{hive_metastore}. Điểm quan trọng là cả Spark (ghi) và Trino (đọc) cùng đăng ký với \emph{một} metastore, nên cả hai thấy cùng một bảng với cùng một snapshot. Không có catalog dùng chung, hai engine sẽ thấy hai thế giới khác nhau.

\subsection{Trino làm engine truy vấn phân tích}
\label{sec:2-trino}

Trino là engine truy vấn SQL phân tán, không tự lưu dữ liệu, mà đọc qua các \emph{connector} \cite{trino_docs}. Connector Iceberg đọc siêu dữ liệu Iceberg qua catalog và đọc tệp Parquet từ HDFS. Có ba lý do Trino được chọn làm đường đọc phân tích thay vì để API gọi Spark:

\begin{itemize}
  \item Trino khởi động truy vấn trong mili-giây tới giây, trong khi một job Spark cần khởi tạo ngữ cảnh. Với bảng điều khiển tương tác, độ trễ này quan trọng.
  \item Trino có sẵn máy khách Python (\code{trino.dbapi}) với tham số hóa truy vấn, phù hợp với lớp FastAPI.
  \item Trino đọc được bảng Iceberg \emph{do Spark ghi} mà không cần sao chép dữ liệu, nên không có hệ thống thứ hai phải đồng bộ.
\end{itemize}

Đoạn mã~\ref{lst:2-trino-catalog} là toàn bộ cấu hình catalog của Trino trong đề tài (đã lược chú thích).

\begin{lstlisting}[caption={Catalog Iceberg của Trino (\texttt{trino/catalog/lakehouse.properties})}, label={lst:2-trino-catalog}]
connector.name=iceberg
iceberg.catalog.type=hive_metastore
hive.metastore.uri=thrift://hive-metastore:9083
hive.config.resources=/etc/hadoop/conf/core-site.xml,/etc/hadoop/conf/hdfs-site.xml
iceberg.file-format=PARQUET
iceberg.compression-codec=SNAPPY
\end{lstlisting}

Trino mở rộng SQL để dùng các khả năng của Iceberg. Đoạn mã~\ref{lst:2-trino-timetravel} liệt kê các snapshot của bảng bronze bằng bảng siêu dữ liệu \texttt{\$snapshots}, rồi đọc bảng ở trạng thái của một snapshot hoặc một thời điểm.

\begin{lstlisting}[language=SQL, caption={Du hành thời gian trong Trino (cú pháp của connector Iceberg)}, label={lst:2-trino-timetravel}]
-- cac snapshot cua bang va thoi diem tao
SELECT snapshot_id, committed_at, operation
FROM lakehouse.bronze."sensor_readings$snapshots"
ORDER BY committed_at DESC;

-- doc bang o trang thai cua mot snapshot
SELECT count(*) FROM lakehouse.bronze.sensor_readings
FOR VERSION AS OF 1234567890123456789;

-- doc bang tai mot thoi diem
SELECT count(*) FROM lakehouse.bronze.sensor_readings
FOR TIMESTAMP AS OF TIMESTAMP '2026-11-01 00:00:00 UTC';
\end{lstlisting}

\paragraph{Giới hạn.} Trino không thay thế được kho gần thời gian thực. Mỗi truy vấn phải đọc siêu dữ liệu Iceberg và tệp Parquet, nên độ trễ nằm ở cỡ trăm mili-giây tới giây, so với cỡ mười mili-giây của tra cứu theo khóa trong Cassandra. Đó là lý do hệ thống có hai đường phục vụ và không trộn nguồn trong cùng một biểu đồ; benchmark ở Mục~\ref{sec:5-bench-lookups} đo chênh lệch này.

